\section{Hoeffding's inequality}
\label{sec:hoeffding-theorem}

We now assemble the main result. The proof is the canonical Chernoff argument:
exponentiate, apply Markov's inequality, factorise the moment generating
function across the independent summands, bound each factor with
\Cref{lem:hoeffding}, optimise the free parameter $\lambda$, and finally combine
the two tails by a union bound.

\begin{theorem}[Hoeffding's inequality]\label{thm:hoeffding}
Let $X_1, \ldots, X_n$ be independent random variables with $X_i \in [a_i, b_i]$
almost surely for each $i$, where $a_i < b_i$, and set $S_n := \sum_{i=1}^n X_i$.
Write $V := \sum_{i=1}^n (b_i - a_i)^2$. Then for every $t > 0$,
\begin{align}\label{eq:hoeffding-main}
\Pr\!\big[\, \abs{S_n - \E[S_n]} \ge t \,\big]
\;\le\; 2 \exp\!\left( - \frac{2 t^2}{\sum_{i=1}^n (b_i - a_i)^2} \right).
\end{align}
\end{theorem}

\begin{proof}
Set $Y_i := X_i - \E[X_i]$, so that $Y_i$ is centered with $Y_i \in [a_i', b_i']$
almost surely for $a_i' := a_i - \E[X_i]$ and $b_i' := b_i - \E[X_i]$; note
$b_i' - a_i' = b_i - a_i$. The centered sum is $S_n - \E[S_n] = \sum_{i=1}^n Y_i$.
The proof proceeds in three stages: a one-sided exponential tail bound, its
optimisation over $\lambda$, and the two-sided combination.

\smallskip\noindent\textbf{Step 1 (Chernoff bound on the upper tail).}
For any $\lambda > 0$, the map $x \mapsto e^{\lambda x}$ is increasing and
positive, so by Markov's inequality applied to $e^{\lambda(S_n - \E S_n)}$,
\begin{align}
\Pr\!\big[\, S_n - \E[S_n] \ge t \,\big]
&\;=\; \Pr\!\big[\, e^{\lambda (S_n - \E S_n)} \ge e^{\lambda t} \,\big] \notag \\
&\;\le\; e^{-\lambda t} \, \E\!\left[ e^{\lambda \sum_{i=1}^n Y_i} \right]
   \;=\; e^{-\lambda t} \prod_{i=1}^n \E\!\left[ e^{\lambda Y_i} \right], \label{eq:chernoff}
\end{align}
where the first step exponentiates both sides of the event with the increasing
bijection $x \mapsto e^{\lambda x}$, the second is Markov's inequality
$\Pr[Z \ge s] \le \E[Z]/s$ applied to the nonnegative variable
$Z = e^{\lambda(S_n - \E S_n)}$, and the last step factorises the expectation of
a product using the independence of $Y_1, \ldots, Y_n$ (a measurable function of
independent variables).

\smallskip\noindent\textbf{Step 2 (apply Hoeffding's lemma and optimise $\lambda$).}
Each $Y_i$ is centered and bounded in an interval of width $b_i - a_i$, so
\Cref{lem:hoeffding} bounds each factor by
$\E[e^{\lambda Y_i}] \le \exp(\lambda^2 (b_i - a_i)^2 / 8)$. Substituting into
Eq.~\eqref{eq:chernoff} and writing $V = \sum_i (b_i - a_i)^2$,
\begin{align}
\Pr\!\big[\, S_n - \E[S_n] \ge t \,\big]
&\;\le\; e^{-\lambda t} \prod_{i=1}^n \exp\!\left( \frac{\lambda^2 (b_i - a_i)^2}{8} \right) \notag \\
&\;=\; \exp\!\left( -\lambda t + \frac{\lambda^2 V}{8} \right)
   \;=:\; \exp\big( g(\lambda) \big), \label{eq:exponent}
\end{align}
where the first step applies \Cref{lem:hoeffding} to each of the $n$ independent
factors, and the second collects the product of exponentials into a single
exponent $g(\lambda) := -\lambda t + \lambda^2 V / 8$. The bound holds for every
$\lambda > 0$, so we minimise $g$. The first-order condition is
\begin{align}
g'(\lambda) \;=\; -t + \frac{\lambda V}{4} \;=\; 0
\qquad\Longleftrightarrow\qquad
\lambda^\star \;=\; \frac{4 t}{V} \;>\; 0, \label{eq:foc}
\end{align}
and $g''(\lambda) = V/4 > 0$ confirms $\lambda^\star$ is the global minimiser.
Evaluating the exponent at $\lambda^\star$,
\begin{align}
g(\lambda^\star)
\;=\; -\frac{4t}{V}\, t + \frac{1}{8}\left( \frac{4t}{V} \right)^{\!2} V
\;=\; -\frac{4 t^2}{V} + \frac{2 t^2}{V}
\;=\; -\frac{2 t^2}{V}, \label{eq:exponent-opt}
\end{align}
where the first equality substitutes $\lambda^\star = 4t/V$ into
$g(\lambda) = -\lambda t + \lambda^2 V/8$, and the remaining equalities are
arithmetic simplification. Combining Eq.~\eqref{eq:exponent}
and Eq.~\eqref{eq:exponent-opt} gives the one-sided bound
\begin{align}
\Pr\!\big[\, S_n - \E[S_n] \ge t \,\big]
\;\le\; \exp\!\left( - \frac{2 t^2}{V} \right). \label{eq:one-sided}
\end{align}

\smallskip\noindent\textbf{Step 3 (two-sided bound by a union bound).}
Applying the one-sided bound Eq.~\eqref{eq:one-sided} to the variables $-X_i \in
[-b_i, -a_i]$, whose interval widths are again $b_i - a_i$ and whose centered sum
is $-(S_n - \E S_n)$, yields the matching lower-tail bound
\begin{align}
\Pr\!\big[\, S_n - \E[S_n] \le -t \,\big]
\;=\; \Pr\!\big[\, (\E S_n - S_n) \ge t \,\big]
\;\le\; \exp\!\left( - \frac{2 t^2}{V} \right). \label{eq:lower-sided}
\end{align}
The two-sided event decomposes as $\{\abs{S_n - \E S_n} \ge t\} = \{S_n - \E S_n
\ge t\} \cup \{S_n - \E S_n \le -t\}$. By a union bound over these two events,
together with Eq.~\eqref{eq:one-sided} and Eq.~\eqref{eq:lower-sided},
\begin{align}
\Pr\!\big[\, \abs{S_n - \E[S_n]} \ge t \,\big]
&\;\le\; \Pr\!\big[\, S_n - \E S_n \ge t \,\big] + \Pr\!\big[\, S_n - \E S_n \le -t \,\big] \notag \\
&\;\le\; 2 \exp\!\left( - \frac{2 t^2}{V} \right), \label{eq:two-sided}
\end{align}
where the first step is the union bound applied to the two tail events, and the
second substitutes the upper- and lower-tail bounds. Since $V = \sum_{i=1}^n
(b_i - a_i)^2$, this is exactly Eq.~\eqref{eq:hoeffding-main}, which completes
the proof.
\end{proof}
